% =============================================================================
% Chapter 3: Testbench Architecture
% =============================================================================
\chapter{Testbench Architecture}
\label{ch:tb_arch}

\section{Testbench Overview}

The testbench file \texttt{spi\_master\_tb\_v2.sv} (module name \texttt{tb\_apb\_spi\_master})
implements a self-checking, directed-test environment for the APB SPI Master IP. It is designed
for the Cadence Xcelium simulator and uses Xcelium-specific waveform dump commands
(\texttt{\$recordfile}, \texttt{\$recordvars}).

The testbench is a single-module, non-UVM environment. This approach was chosen for simplicity,
fast simulation, and direct correlation between stimulus and DUT behavior. The same architectural
pattern can be reused for other APB peripherals with minimal changes (see Chapter~\ref{ch:apb_guide}).

\section{Testbench Parameters}

\begin{table}[H]
\centering
\caption{Testbench Parameters}
\label{tab:tb_params}
\begin{tabularx}{\textwidth}{@{}lcX@{}}
\toprule
\textbf{Parameter} & \textbf{Value} & \textbf{Description} \\
\midrule
\texttt{BUFFER\_DEPTH}   & 10    & FIFO depth, matches DUT parameter. \\
\texttt{APB\_ADDR\_WIDTH} & 12   & APB address width in bits (4\,KB space). \\
\texttt{CLK\_PERIOD}      & 10\,ns & System clock period (100\,MHz). \\
\bottomrule
\end{tabularx}
\end{table}

\section{Clock and Reset Generation}

The system clock (\texttt{HCLK}) is generated using the standard Verilog idiom:
\begin{lstlisting}
initial HCLK = 0;
always #(CLK_PERIOD/2) HCLK = ~HCLK;  // 100 MHz
\end{lstlisting}

Reset (\texttt{HRESETn}) is active-low. The testbench asserts reset for 5 clock cycles at
simulation start, then deasserts it and waits 3 additional cycles before running testcases.
TC14 re-asserts reset mid-simulation (8 cycles) to test post-reset interrupt behavior.

\section{APB Bus Functional Model (BFM)}
\label{sec:apb_bfm}

The testbench implements APB read and write transactions as SystemVerilog tasks. These tasks
follow the AMBA 3 APB protocol with two-cycle transfers (setup + access):

\subsection{APB Write Task}

\begin{lstlisting}
task automatic apb_write(input logic [APB_ADDR_WIDTH-1:0] addr,
                         input logic [31:0]               data);
    @(posedge HCLK); #1;              // Setup phase
    PADDR = addr; PWDATA = data;
    PWRITE = 1'b1; PSEL = 1'b1; PENABLE = 1'b0;
    @(posedge HCLK); #1;              // Access phase
    PENABLE = 1'b1;
    @(posedge HCLK); #1;              // Deassert
    PSEL = 1'b0; PENABLE = 1'b0; PWRITE = 1'b0;
endtask
\end{lstlisting}

The \texttt{\#1} delays after each \texttt{@(posedge HCLK)} ensure that signal changes occur in
the \emph{postponed region} of the simulation cycle, avoiding race conditions with the DUT's
registered logic. This is a standard best practice for APB testbenches.

\subsection{APB Read Task}

\begin{lstlisting}
task automatic apb_read(input logic [APB_ADDR_WIDTH-1:0] addr,
                        output logic [31:0]              rdata);
    @(posedge HCLK); #1;              // Setup phase
    PADDR = addr; PWRITE = 1'b0;
    PSEL = 1'b1; PENABLE = 1'b0;
    @(posedge HCLK); #1;              // Access phase
    PENABLE = 1'b1;
    @(posedge HCLK);                  // Sample BEFORE #1
    rdata = PRDATA;                   // Capture at posedge
    #1;
    PSEL = 1'b0; PENABLE = 1'b0;
endtask
\end{lstlisting}

\textbf{Critical timing detail:} For RXFIFO reads, \texttt{PRDATA} must be sampled at the
\texttt{posedge HCLK} \emph{before} the \texttt{\#1} delay. This is because the FIFO's
\texttt{pointer\_out} increments on the same clock edge (via non-blocking assignment). If we
sampled after \texttt{\#1}, the pointer would have already advanced and the combinational
\texttt{data\_o = buffer[pointer\_out]} would show the \emph{next} word. This subtlety was
discovered during TC4/TC10 debugging.

\subsection{Address Calculation}

The register address function translates a 4-bit register index to a 12-bit APB address:
\begin{lstlisting}
function automatic logic [APB_ADDR_WIDTH-1:0] reg_addr(input logic [3:0] idx);
    return {6'h0, idx, 2'b00};  // Word-aligned: index x 4
endfunction
\end{lstlisting}

\section{SPI Slave BFM}
\label{sec:slave_bfm}

The testbench includes a hardware-style SPI slave BFM to drive data onto the \texttt{spi\_sdi}
lines during read transactions. This BFM evolved through several iterations to solve race
condition issues specific to Xcelium (see Chapter~\ref{ch:debugging}).

\subsection{Design Rationale}

The DUT's \texttt{spi\_clk} output is a \emph{registered} version of the internal clock, meaning
it lags the internal \texttt{spi\_rise} event by one system clock cycle. A task-based slave BFM
that clocks on \texttt{@(posedge spi\_clk)} would change data too late---the DUT's RX register
would sample the old data. The solution is a counter-based, always-block BFM:

\begin{lstlisting}
logic        slave_drive_en;
logic [31:0] slave_tx_data;
logic  [4:0] slave_bit_cnt;

always_ff @(negedge spi_clk or negedge HRESETn) begin
    if (!HRESETn)          slave_bit_cnt <= 5'd0;
    else if (slave_drive_en) slave_bit_cnt <= slave_bit_cnt + 5'd1;
    else                     slave_bit_cnt <= 5'd0;
end

wire spi_sdi0 = spi_sdo0;  // Loopback
wire spi_sdi1 = slave_drive_en ?
     slave_tx_data[5'd31 - slave_bit_cnt] : spi_sdo1;
wire spi_sdi2 = spi_sdo2;  // Loopback
wire spi_sdi3 = spi_sdo3;  // Loopback
\end{lstlisting}

\subsection{Key Points}

\begin{itemize}[leftmargin=2em]
  \item \textbf{Wire types:} The \texttt{spi\_sdi*} signals are declared as \texttt{wire}, not
        \texttt{logic}. In Xcelium, using \texttt{assign} on a \texttt{logic} type can silently
        fail, leaving the signal stuck at zero. Using \texttt{wire} avoids this simulator-specific
        pitfall.
  \item \textbf{Negedge clocking:} The bit counter increments on the \emph{falling edge} of
        \texttt{spi\_clk}. Since the DUT samples on the rising edge, the slave data is guaranteed
        stable before the next sample point.
  \item \textbf{MSB-first:} Data is indexed as \texttt{slave\_tx\_data[31 - slave\_bit\_cnt]},
        matching the DUT's MSB-first shift register.
  \item \textbf{Loopback on sdi0, sdi2, sdi3:} When the slave BFM is not driving, these lines
        loop back from the corresponding \texttt{sdo} lines for basic connectivity verification.
\end{itemize}

\section{EOT Flag and Handling}

The testbench latches the EOT event using a sticky flag:
\begin{lstlisting}
logic eot_flag;
always_ff @(posedge HCLK or negedge HRESETn) begin
    if (!HRESETn)       eot_flag <= 1'b0;
    else if (events_o[1]) eot_flag <= 1'b1;
end
\end{lstlisting}

A \texttt{clear\_eot\_flag} task resets this flag before each testcase using \texttt{force}/\texttt{release}:
\begin{lstlisting}
task clear_eot_flag;
    @(negedge HCLK);
    force eot_flag = 1'b0;
    @(negedge HCLK);
    release eot_flag;
endtask
\end{lstlisting}

The \texttt{wait\_eot} task polls \texttt{eot\_flag} with a configurable timeout (default 20,000 cycles)
to detect transaction completion.

\section{Pass/Fail Checking Infrastructure}

The \texttt{check} task implements masked comparison with formatted output:
\begin{lstlisting}
task automatic check(input string msg,
                     input logic [31:0] got,
                     input logic [31:0] exp,
                     input logic [31:0] mask = 32'hFFFF_FFFF);
    if ((got & mask) === (exp & mask)) begin
        $display("  [PASS] %s  got=0x%08X  exp=0x%08X", msg, got, exp);
        pass_cnt++;
    end else begin
        $display("  [FAIL] %s  got=0x%08X  exp=0x%08X (mask=0x%08X)",
                 msg, got, exp, mask);
        fail_cnt++;
    end
endtask
\end{lstlisting}

Key features:
\begin{itemize}[leftmargin=2em]
  \item Uses \texttt{===} (four-state equality) to correctly handle X/Z values.
  \item Supports bit masking for partial register checks (e.g., checking only status bits [6:0]).
  \item Maintains running \texttt{pass\_cnt} and \texttt{fail\_cnt} counters for the final summary.
  \item Provides self-documenting output with both actual and expected values.
\end{itemize}

\section{Helper Tasks}

\subsection{\texttt{cfg\_spi} -- SPI Configuration}

Writes all five configuration registers (CLKDIV, SPICMD, SPIADR, SPILEN, SPIDUM) in a single
call, simplifying testcase setup.

\subsection{\texttt{trigger} -- Transaction Start}

Writes the STATUS register with the operation code (rd/wr/qrd/qwr) and chip-select value to
start an SPI transaction.

\subsection{\texttt{push\_tx} / \texttt{pop\_rx} -- FIFO Access}

Simple wrappers around \texttt{apb\_write(REG\_TXFIFO)} and \texttt{apb\_read(REG\_RXFIFO)}.

\section{Global Watchdog}

A global watchdog timer terminates simulation after 10\,ms (10,000,000~ns) to prevent infinite
hangs due to unexpected FSM lockups:
\begin{lstlisting}
initial begin
    #10_000_000;
    $display("[ERROR] Global watchdog timeout!");
    $finish;
end
\end{lstlisting}

\section{Waveform Dumping}

For Cadence Xcelium/SimVision, the testbench uses:
\begin{lstlisting}
initial begin
    $recordfile("waves_v2.trn");
    $recordvars();
end
\end{lstlisting}

This records all signals in the design hierarchy to a \texttt{.trn} file viewable in SimVision.
For other simulators, replace with \texttt{\$dumpfile}/\texttt{\$dumpvars} (VCD) or simulator-specific
commands.
